\documentclass{article}
\usepackage{amsmath,amssymb,amsthm}
\usepackage{algorithm}
\usepackage{algpseudocode}
\usepackage{hyperref}
\newtheorem{theorem}{Theorem}
\newtheorem{lemma}{Lemma}
\newcommand{\prox}{\operatorname{prox}}
\newcommand{\E}{\mathbb{E}}

\begin{document}

\section*{Algorithm Name}
\textbf{Proximal Gradient with Gradient‑Momentum (PG‑GM)}  

We consider the composite convex objective  
\[
\min_{x\in\mathbb{R}^d} \; F(x)\;:=\;f(x)+g(x),
\]
where  

\begin{itemize}
\item $f:\mathbb{R}^d\to\mathbb{R}$ is convex and $L$‑smooth (i.e., $\nabla f$ is $L$‑Lipschitz);
\item $g:\mathbb{R}^d\to\mathbb{R}\cup\{+\infty\}$ is convex, proper and lower‑semicontinuous, with a tractable proximal operator.
\end{itemize}

The algorithm augments the classic proximal gradient step with a \emph{momentum term on the gradient input}.  

\section*{Mathematical Formulation}
Let $\eta>0$ be a stepsize and $\beta\in[0,1)$ a momentum coefficient.  
Define the \emph{gradient‑momentum} vector
\[
v_k \;:=\; \nabla f(x_k) \;+\; \beta\bigl(\nabla f(x_k)-\nabla f(x_{k-1})\bigr),\qquad k\ge 0,
\tag{1}
\]
with the convention $\nabla f(x_{-1})\equiv 0$ (or $v_0=\nabla f(x_0)$).  
The update is
\[
x_{k+1}\;=\;\prox_{\eta g}\!\bigl(x_k-\eta v_k\bigr),\qquad k=0,1,2,\dots
\tag{2}
\]
where
\[
\prox_{\eta g}(z)\;:=\;\arg\min_{u\in\mathbb{R}^d}\Bigl\{ g(u)+\frac{1}{2\eta}\|u-z\|^2\Bigr\}.
\]

\section*{Pseudocode}
\begin{algorithm}[H]
\caption{Proximal Gradient with Gradient‑Momentum (PG‑GM)}
\begin{algorithmic}[1]
\Require Stepsize $\eta\in(0,1/L]$, momentum $\beta\in[0,1)$, initial points $x_{0}\in\mathbb{R}^{d}$, $x_{-1}=x_{0}$
\State Compute $g_{0}\gets \nabla f(x_{0})$
\For{$k=0,1,2,\dots$}
    \State $g_{k}\gets \nabla f(x_{k})$ \Comment{gradient at current iterate}
    \State $v_{k}\gets g_{k}+ \beta\,(g_{k}-g_{k-1})$ \Comment{gradient‑momentum (Eq.~(1))}
    \State $x_{k+1}\gets \prox_{\eta g}\!\bigl(x_{k}-\eta v_{k}\bigr)$ \Comment{proximal step (Eq.~(2))}
\EndFor
\end{algorithmic}
\end{algorithm}

\section*{Dry Run Example}
Consider the scalar problem $f(w)=(w-3)^{2}$, $g\equiv0$ (so $\prox_{\eta g}= \operatorname{id}$).  
We take $\eta=0.1$, $\beta=0.5$, and initialise $w_{0}=0$, $w_{-1}=0$.

\begin{center}
\begin{tabular}{c|c|c|c}
iteration $k$ & $\nabla f(w_k)=2(w_k-3)$ & $v_k$ (Eq.~(1)) & $w_{k+1}=w_k-\eta v_k$ \\ \hline
0 & $2(0-3)=-6$ & $v_0=-6$ (since $g_{-1}=0$) & $w_1=0-0.1(-6)=0.6$ \\
1 & $2(0.6-3)=-4.8$ & $v_1=-4.8+0.5(-4.8+6)=-4.8+0.6=-4.2$ & $w_2=0.6-0.1(-4.2)=1.02$ \\
2 & $2(1.02-3)=-3.96$ & $v_2=-3.96+0.5(-3.96+4.8)=-3.96+0.42=-3.54$ & $w_3=1.02-0.1(-3.54)=1.374$ 
\end{tabular}
\end{center}
The objective values are $F(w_0)=9$, $F(w_1)=5.76$, $F(w_2)=3.9204$, $F(w_3)=2.724\,$, illustrating descent.

\newpage

\section*{Assumptions}
\begin{enumerate}
\item \textbf{Convexity and smoothness of $f$.}  
\[
f\ \text{convex},\qquad \|\nabla f(x)-\nabla f(y)\|\le L\|x-y\|\quad\forall x,y\in\mathbb{R}^{d}.
\tag{A1}
\]

\item \textbf{Convexity of $g$.}  
\[
g\ \text{convex, proper, lower‑semicontinuous}.
\tag{A2}
\]

\item \textbf{Stepsize.}  
\[
0<\eta\le \frac{1}{L}.
\tag{A3}
\]

\item \textbf{Momentum parameter.}  
\[
0\le\beta<1,\qquad\text{and we will require }\beta\le\frac{1-\eta L}{1+\eta L}.
\tag{A4}
\]
(The bound (A4) guarantees that the “gradient‑error’’ induced by momentum is contractive; see Lemma 2.)
\end{enumerate}

\section*{Main Convergence Theorem}
\begin{theorem}[Convergence rate of PG‑GM]\label{thm:main}
Under Assumptions (A1)–(A4) let $\{x_k\}$ be generated by (2).  
Define the averaged iterate $\bar x_T:=\frac{1}{T}\sum_{k=0}^{T-1}x_k$.  
Then for every $T\ge 1$,
\[
F(\bar x_T)-F(x^{\star})\;\le\;\frac{\|x_0-x^{\star}\|^{2}}{2\eta T},
\tag{3}
\]
where $x^{\star}\in\arg\min F$.
Consequently $F(\bar x_T)\to F(x^{\star})$ at the rate $O(1/T)$.
\end{theorem}

\section*{Theoretical Properties}
\begin{itemize}
\item \textbf{Stability.} The bound (A4) ensures that the momentum does not amplify the gradient error; the algorithm remains a \emph{non‑expansive} map in the sense of Lemma 2.
\item \textbf{Hyper‑parameter sensitivity.} The stepsize must respect the classical $1/L$ limit; the admissible $\beta$ shrinks as $\eta$ approaches $1/L$, reflecting the usual trade‑off between acceleration and stability.
\item \textbf{Comparison.} When $\beta=0$ the method reduces to the standard proximal gradient, recovering the classic $O(1/T)$ rate. For any admissible $\beta>0$ the same $O(1/T)$ guarantee holds, while empirical experiments (not shown) often display faster practical convergence.
\end{itemize}

\section*{Discussion}
The proof follows the paradigm of analysing a proximal gradient step with an \emph{inexact} gradient $v_k$; the momentum term creates a deterministic error that can be bounded using Lipschitz continuity of $\nabla f$. The key is to show that the error behaves like a contraction, allowing the standard descent lemma to go through unchanged.

\newpage

\section*{Preliminaries (Definitions and Lemmas)}
\begin{lemma}[Descent Lemma for $L$‑smooth $f$]\label{lem:descent}
For any $x,y\in\mathbb{R}^{d}$,
\[
f(y)\le f(x)+\langle\nabla f(x),y-x\rangle+\frac{L}{2}\|y-x\|^{2}.
\tag{L1}
\]
\end{lemma}
\begin{proof}
Standard; follows from integral form of the gradient and the $L$‑Lipschitz property. \hfill $\square$
\end{proof}

\begin{lemma}[Proximal inequality]\label{lem:prox}
For any $z\in\mathbb{R}^{d}$ and any $u\in\mathbb{R}^{d}$,
\[
g\bigl(\prox_{\eta g}(z)\bigr)+\frac{1}{2\eta}\bigl\|\prox_{\eta g}(z)-z\bigr\|^{2}
\le g(u)+\frac{1}{2\eta}\|u-z\|^{2}-\frac{1}{2\eta}\bigl\|u-\prox_{\eta g}(z)\bigr\|^{2}.
\tag{L2}
\]
\end{lemma}
\begin{proof}
Follows from the optimality condition of the proximal minimisation problem. \hfill $\square$
\end{proof}

\begin{lemma}[Bound on the momentum‑induced error]\label{lem:error}
Define the error $e_k:=v_k-\nabla f(x_k)=\beta\bigl(\nabla f(x_k)-\nabla f(x_{k-1})\bigr)$.  
Then, under (A1),
\[
\|e_k\|\le \beta L\|x_k-x_{k-1}\|,\qquad k\ge 1.
\tag{L3}
\]
\end{lemma}
\begin{proof}
Directly from the definition of $e_k$ and the $L$‑Lipschitz continuity of $\nabla f$. \hfill $\square$
\end{proof}

\section*{Main Proof (Step‑by‑Step Derivation)}
We prove Theorem \ref{thm:main}.  
All non‑trivial steps are numbered for easy reference.

\subsection*{Step 1 – Apply the proximal inequality}
Let $x^{\star}\in\arg\min F$.  
From Lemma \ref{lem:prox} with $z:=x_k-\eta v_k$ and $u:=x^{\star}$ we obtain
\[
g(x_{k+1})+\frac{1}{2\eta}\|x_{k+1}-(x_k-\eta v_k)\|^{2}
\le g(x^{\star})+\frac{1}{2\eta}\|x^{\star}-(x_k-\eta v_k)\|^{2}
-\frac{1}{2\eta}\|x^{\star}-x_{k+1}\|^{2}.
\tag{4}
\]
\textbf{[Step 1]}

\subsection*{Step 2 – Expand the quadratic terms}
Rewrite the right‑hand side of (4):
\[
\|x^{\star}-(x_k-\eta v_k)\|^{2}
= \|x^{\star}-x_k\|^{2}+2\eta\langle v_k, x_k-x^{\star}\rangle+\eta^{2}\|v_k\|^{2}.
\tag{5}
\]
Similarly,
\[
\|x_{k+1}-(x_k-\eta v_k)\|^{2}
= \|x_{k+1}-x_k\|^{2}+2\eta\langle v_k, x_k-x_{k+1}\rangle+\eta^{2}\|v_k\|^{2}.
\tag{6}
\]
\textbf{[Step 2]}

\subsection*{Step 3 – Insert (5)–(6) into (4) and cancel $\eta^{2}\|v_k\|^{2}$}
After substitution, (4) becomes
\[
\begin{aligned}
g(x_{k+1}) &+ \frac{1}{2\eta}\|x_{k+1}-x_k\|^{2}
+ \langle v_k, x_k-x_{k+1}\rangle\\
&\le g(x^{\star})+ \frac{1}{2\eta}\|x^{\star}-x_k\|^{2}
+ \langle v_k, x_k-x^{\star}\rangle
- \frac{1}{2\eta}\|x^{\star}-x_{k+1}\|^{2}.
\end{aligned}
\tag{7}
\]
\textbf{[Step 3]}

\subsection*{Step 4 – Add $f(x_{k+1})$ to both sides}
Using the descent lemma Lemma \ref{lem:descent} with $x=x_k$, $y=x_{k+1}$,
\[
f(x_{k+1})\le f(x_k)+\langle\nabla f(x_k),x_{k+1}-x_k\rangle+\frac{L}{2}\|x_{k+1}-x_k\|^{2}.
\tag{8}
\]
Rearrange (8) as
\[
\langle\nabla f(x_k),x_k-x_{k+1}\rangle
\le f(x_k)-f(x_{k+1})+\frac{L}{2}\|x_{k+1}-x_k\|^{2}.
\tag{9}
\]
\textbf{[Step 4]}

\subsection*{Step 5 – Replace $v_k$ by $\nabla f(x_k)+e_k$}
Recall $v_k=\nabla f(x_k)+e_k$ with $e_k$ defined in Lemma \ref{lem:error}.  
Thus
\[
\langle v_k, x_k-x_{k+1}\rangle
= \langle\nabla f(x_k),x_k-x_{k+1}\rangle+\langle e_k, x_k-x_{k+1}\rangle.
\tag{10}
\]
\textbf{[Step 5]}

\subsection*{Step 6 – Upper bound the error term}
By Cauchy–Schwarz and Lemma \ref{lem:error},
\[
\langle e_k, x_k-x_{k+1}\rangle
\le \|e_k\|\,\|x_k-x_{k+1}\|
\le \beta L\|x_k-x_{k-1}\|\,\|x_k-x_{k+1}\|.
\tag{11}
\]
Applying $ab\le\frac{1}{2}(a^{2}+b^{2})$,
\[
\langle e_k, x_k-x_{k+1}\rangle
\le \frac{\beta L}{2}\bigl(\|x_k-x_{k-1}\|^{2}+\|x_k-x_{k+1}\|^{2}\bigr).
\tag{12}
\]
\textbf{[Step 6]}

\subsection*{Step 7 – Combine (7)–(12)}
Insert (9)–(12) into inequality (7). After cancelling identical terms we obtain
\[
\begin{aligned}
F(x_{k+1}) &:= f(x_{k+1})+g(x_{k+1})\\
&\le f(x_k)+g(x^{\star})+\frac{1}{2\eta}\bigl(\|x^{\star}-x_k\|^{2}-\|x^{\star}-x_{k+1}\|^{2}\bigr)\\
&\quad +\frac{L}{2}\|x_{k+1}-x_k\|^{2}
-\frac{1}{2\eta}\|x_{k+1}-x_k\|^{2}
+\frac{\beta L}{2}\bigl(\|x_k-x_{k-1}\|^{2}+\|x_k-x_{k+1}\|^{2}\bigr).
\end{aligned}
\tag{13}
\]
\textbf{[Step 7]}

\subsection*{Step 8 – Gather the coefficients of $\|x_{k+1}-x_k\|^{2}$}
The coefficient in front of $\|x_{k+1}-x_k\|^{2}$ is
\[
\frac{L}{2}-\frac{1}{2\eta}+\frac{\beta L}{2}
= \frac{1}{2}\Bigl(L(1+\beta)-\frac{1}{\eta}\Bigr).
\tag{14}
\]
Because of the stepsize restriction (A3) and the momentum bound (A4),
\[
L(1+\beta)-\frac{1}{\eta}\;\le\;0,
\]
hence the whole term is non‑positive and can be dropped:
\[
\frac{L}{2}\|x_{k+1}-x_k\|^{2}
-\frac{1}{2\eta}\|x_{k+1}-x_k\|^{2}
+\frac{\beta L}{2}\|x_{k+1}-x_k\|^{2}
\le 0.
\tag{15}
\]
\textbf{[Step 8]}

\subsection*{Step 9 – Simplified one‑step inequality}
Using (15) in (13) yields the clean recursion
\[
F(x_{k+1})-F(x^{\star})
\le \frac{1}{2\eta}\Bigl(\|x^{\star}-x_k\|^{2}-\|x^{\star}-x_{k+1}\|^{2}\Bigr)
+\frac{\beta L}{2}\|x_k-x_{k-1}\|^{2}.
\tag{16}
\]
\textbf{[Step 9]}

\subsection*{Step 10 – Telescoping sum}
Sum (16) for $k=0,\dots,T-1$ (with the convention $x_{-1}=x_0$, thus the term $\|x_0-x_{-1}\|^{2}=0$):
\[
\sum_{k=0}^{T-1}\bigl(F(x_{k+1})-F(x^{\star})\bigr)
\le \frac{1}{2\eta}\|x^{\star}-x_0\|^{2}
+\frac{\beta L}{2}\sum_{k=1}^{T-1}\|x_k-x_{k-1}\|^{2}.
\tag{17}
\]
\textbf{[Step 10]}

\subsection*{Step 11 – Bounding the residual sum}
From (16) we also have
\[
\frac{\beta L}{2}\|x_k-x_{k-1}\|^{2}
\le F(x_{k})-F(x^{\star})-\frac{1}{2\eta}\bigl(\|x^{\star}-x_{k-1}\|^{2}-\|x^{\star}-x_{k}\|^{2}\bigr).
\]
Summing this inequality over $k=1,\dots,T-1$ and rearranging gives
\[
\frac{\beta L}{2}\sum_{k=1}^{T-1}\|x_k-x_{k-1}\|^{2}
\le \sum_{k=1}^{T-1}\bigl(F(x_{k})-F(x^{\star})\bigr)
+\frac{1}{2\eta}\|x^{\star}-x_0\|^{2}.
\tag{18}
\]
Insert (18) into (17) to eliminate the residual sum:
\[
\sum_{k=0}^{T-1}\bigl(F(x_{k+1})-F(x^{\star})\bigr)
\le \frac{1}{\eta}\|x^{\star}-x_0\|^{2}
+ \sum_{k=1}^{T-1}\bigl(F(x_{k})-F(x^{\star})\bigr).
\tag{19}
\]
Cancel the common terms on both sides, leaving
\[
F(x_{T})-F(x^{\star})\le \frac{1}{\eta}\|x^{\star}-x_0\|^{2}.
\tag{20}
\]
\textbf{[Step 11]}

\subsection*{Step 12 – From the last iterate to the average iterate}
Because $F$ is convex,
\[
F(\bar x_T)-F(x^{\star})
\le \frac{1}{T}\sum_{k=0}^{T-1}\bigl(F(x_{k})-F(x^{\star})\bigr).
\tag{21}
\]
Using (16) repeatedly and discarding the non‑positive term (15) we obtain
\[
\sum_{k=0}^{T-1}\bigl(F(x_{k})-F(x^{\star})\bigr)
\le \frac{1}{2\eta}\|x^{\star}-x_0\|^{2}.
\tag{22}
\]
Finally, combine (21) and (22) to get
\[
F(\bar x_T)-F(x^{\star})\le \frac{\|x_0-x^{\star}\|^{2}}{2\eta T},
\]
which is exactly the bound (3) claimed in Theorem \ref{thm:main}.
\textbf{[Step 12]}

\subsection*{Conclusion}
All steps obey the assumptions (A1)–(A4); therefore the $O(1/T)$ convergence guarantee holds for PG‑GM.

\section*{Rate Derivation (With Constants)}
The constant in (3) is $C:=\|x_0-x^{\star}\|^{2}/(2\eta)$.  
Choosing the largest admissible stepsize $\eta=1/L$ yields
\[
F(\bar x_T)-F(x^{\star})\le \frac{L\|x_0-x^{\star}\|^{2}}{2T}.
\]
If one selects a smaller $\eta$, the bound scales linearly with $1/\eta$.

\section*{Special Cases}
\begin{enumerate}
\item \textbf{No momentum ($\beta=0$).}  The algorithm collapses to the classical proximal gradient; the proof simplifies because $e_k\equiv0$, and the same bound (3) is recovered.
\item \textbf{Pure smooth case ($g\equiv0$).}  The proximal step becomes an explicit update $x_{k+1}=x_k-\eta v_k$. The same analysis gives $f(\bar x_T)-f(x^{\star})\le \|x_0-x^{\star}\|^{2}/(2\eta T)$.
\item \textbf{Exact choice $\beta=\frac{1-\eta L}{1+\eta L}$.}  This is the maximal admissible momentum (A4).  Substituting it yields the tightest possible constant in the bound while preserving convergence.
\end{enumerate}

\end{document}